\documentclass[10pt,a4paper]{amsart}
\usepackage{amsmath,amssymb,amsthm,graphicx}
% Retained minimal preamble for the Stacks Project chapter varieties.tex.
% Verbatim excerpt of the chapter preamble: Xy-pic setup (native preamble lines 14-18), the native theorem
% declarations with the mathematical macros (native preamble lines 154-217), and the project own no-op
% definitions of the reference/slogan/history display environments (native wrapper.tex lines 1-3) that the
% counted statement uses for its slogan. Copied mechanically; nothing added, removed or redefined.
% For commutative diagrams we use Xy-pic
\usepackage[all]{xy}

% We use 2cell for 2-commutative diagrams.
\xyoption{2cell}
% Theorem environments.
%
\theoremstyle{plain}
\newtheorem{theorem}[subsection]{Theorem}
\newtheorem{proposition}[subsection]{Proposition}
\newtheorem{lemma}[subsection]{Lemma}

\theoremstyle{definition}
\newtheorem{definition}[subsection]{Definition}
\newtheorem{example}[subsection]{Example}
\newtheorem{exercise}[subsection]{Exercise}
\newtheorem{situation}[subsection]{Situation}

\theoremstyle{remark}
\newtheorem{remark}[subsection]{Remark}
\newtheorem{remarks}[subsection]{Remarks}

\numberwithin{equation}{subsection}

% Macros
%
\def\lim{\mathop{\mathrm{lim}}\nolimits}
\def\colim{\mathop{\mathrm{colim}}\nolimits}
\def\Spec{\mathop{\mathrm{Spec}}}
\def\Hom{\mathop{\mathrm{Hom}}\nolimits}
\def\Ext{\mathop{\mathrm{Ext}}\nolimits}
\def\SheafHom{\mathop{\mathcal{H}\!\mathit{om}}\nolimits}
\def\SheafExt{\mathop{\mathcal{E}\!\mathit{xt}}\nolimits}
\def\Sch{\mathit{Sch}}
\def\Mor{\mathop{\mathrm{Mor}}\nolimits}
\def\Ob{\mathop{\mathrm{Ob}}\nolimits}
\def\Sh{\mathop{\mathit{Sh}}\nolimits}
\def\NL{\mathop{N\!L}\nolimits}
\def\CH{\mathop{\mathrm{CH}}\nolimits}
\def\proetale{{pro\text{-}\acute{e}tale}}
\def\etale{{\acute{e}tale}}
\def\QCoh{\mathit{QCoh}}
\def\Ker{\mathop{\mathrm{Ker}}}
\def\Im{\mathop{\mathrm{Im}}}
\def\Coker{\mathop{\mathrm{Coker}}}
\def\Coim{\mathop{\mathrm{Coim}}}

% Boxtimes
%
\DeclareMathSymbol{\boxtimes}{\mathbin}{AMSa}{"02}

%
% Macros for moduli stacks/spaces
%
\def\QCohstack{\mathcal{QC}\!\mathit{oh}}
\def\Cohstack{\mathcal{C}\!\mathit{oh}}
\def\Spacesstack{\mathcal{S}\!\mathit{paces}}
\def\Quotfunctor{\mathrm{Quot}}
\def\Hilbfunctor{\mathrm{Hilb}}
\def\Curvesstack{\mathcal{C}\!\mathit{urves}}
\def\Polarizedstack{\mathcal{P}\!\mathit{olarized}}
\def\Complexesstack{\mathcal{C}\!\mathit{omplexes}}
% \Pic is the operator that assigns to X its picard group, usage \Pic(X)
% \Picardstack_{X/B} denotes the Picard stack of X over B
% \Picardfunctor_{X/B} denotes the Picard functor of X over B
\def\Pic{\mathop{\mathrm{Pic}}\nolimits}
\def\Picardstack{\mathcal{P}\!\mathit{ic}}
\def\Picardfunctor{\mathrm{Pic}}
\def\Deformationcategory{\mathcal{D}\!\mathit{ef}}
\newenvironment{reference}{}{}
\newenvironment{slogan}{}{}
\newenvironment{history}{}{}
\usepackage{cleveref}
\crefname{originaldefinition}{original definition}{original definitions}
\Crefname{originaldefinition}{Original definition}{Original definitions}
\makeatletter
\newlabel{algebra-definition-geometrically-reduced}{{030S}{source}{}{sourcepointer-1}{}}
\newlabel{algebra-definition-geometrically-reduced@cref}{{[originaldefinition][0][]030S}{[1][0][]source}{}{}{}}
\makeatother
\crefname{originalsection}{original section}{original sections}
\Crefname{originalsection}{Original section}{Original sections}
\makeatletter
\newlabel{properties-section-integral}{{01OJ}{source}{}{sourcepointer-2}{}}
\newlabel{properties-section-integral@cref}{{[originalsection][0][]01OJ}{[1][0][]source}{}{}{}}
\makeatother
\crefname{originalsection}{original section}{original sections}
\Crefname{originalsection}{Original section}{Original sections}
\makeatletter
\newlabel{schemes-section-fibre-products}{{01JO}{source}{}{sourcepointer-3}{}}
\newlabel{schemes-section-fibre-products@cref}{{[originalsection][0][]01JO}{[1][0][]source}{}{}{}}
\makeatother
\crefname{originallemma}{original lemma}{original lemmas}
\Crefname{originallemma}{Original lemma}{Original lemmas}
\makeatletter
\newlabel{algebra-lemma-geometrically-reduced-finite-purely-inseparable-extension}{{030V}{source}{}{sourcepointer-4}{}}
\newlabel{algebra-lemma-geometrically-reduced-finite-purely-inseparable-extension@cref}{{[originallemma][0][]030V}{[1][0][]source}{}{}{}}
\makeatother
\title{Seven equivalent characterisations of geometric reducedness in positive characteristic}
\author{The Stacks Project Authors}
\date{stacks-project:varieties.tex:local-snapshot-20261011; GNU FDL 1.2 or later}
\begin{document}
\maketitle
\noindent\textit{Editorial note.} All delivered mathematical text is the author's own, copied from the pinned Stacks Project chapter file varieties.tex; no statement or proof step was rewritten, reordered or removed.\par\medskip
\noindent\textit{Editorial note.} The preamble is the chapter's own retained minimal preamble (sources/stacks-varieties/retained-preamble.tex), a verbatim excerpt of the native Xy-pic setup, of the native theorem declarations with mathematical macros, and of the project's own no-op definitions of the reference display environments; no macro was invented or redefined.\par\medskip
\noindent\textit{Editorial note.} All four referenced results live in earlier chapters, so their locators point at the verbatim lines of the project's own tag list staged in the snapshot (sources/stacks-chow/tags.txt) that carry the canonical tags.\par\medskip
\noindent\textit{Editorial note.} Overlap pre-check: this label occurs in no other staged chapter (run with agent-authoring/stk/card\_plan.py varieties before opening).\par\medskip
\noindent\textit{Allowed standard prerequisites (editorial).}\begin{itemize}
\item Geometrically reduced algebra over a field, and reducedness transferred to fibres. Algebra, Definition labelled algebra-definition-geometrically-reduced (canonical Stacks tag 030S); named in the counted statement at line 431. Supplies the algebraic notion used in condition (7); the passage between the ring conditions is the counted content.
\item Reducedness of a ring and of its base change, and integrality of the structure sheaf. Properties, Section labelled properties-section-integral (canonical Stacks tag 01OJ); invoked at proof lines 444-445. Supplies that reducedness of Spec(O(U) tensor\_k k') is reducedness of the ring, and the gluing description of the base change.
\item Fibre products of schemes are obtained by gluing the spectra of the tensor products over an affine cover. Schemes, Section labelled schemes-section-fibre-products (canonical Stacks tag 01JO); invoked at proof lines 450-452. Supplies the description of X\_\{k'\} by gluing spectra of O\_X(U) tensor\_k k', used in the implication (7) => (1).
\item It suffices to test geometric reducedness on finite purely inseparable extensions. Algebra, Lemma labelled algebra-lemma-geometrically-reduced-finite-purely-inseparable-extension (canonical Stacks tag 030V); invoked at proof lines 457-459. Supplies the equivalence of (1) with conditions (3)-(6) by applying it to the rings O\_X(U).
\end{itemize}\medskip
\ifdefined\hypertarget\hypertarget{sourcepointer-1}{}\fi
\noindent\textit{Source reference.} Original definition 030S: sources/stacks-chow/tags.txt, source lines 3959--3959. Cross-chapter pointer named in the counted statement at line 431 for the ring-theoretic notion; the equivalence proof is retained in the counted span.\par\medskip
\ifdefined\hypertarget\hypertarget{sourcepointer-2}{}\fi
\noindent\textit{Source reference.} Original section 01OJ: sources/stacks-chow/tags.txt, source lines 2196--2196. Cross-chapter pointer invoked once at proof lines 444-445 for the reading of reducedness of the spectrum as reducedness of the ring.\par\medskip
\ifdefined\hypertarget\hypertarget{sourcepointer-3}{}\fi
\noindent\textit{Source reference.} Original section 01JO: sources/stacks-chow/tags.txt, source lines 2021--2021. Cross-chapter pointer invoked once at proof lines 450-452 for the gluing description of the base change X\_\{k'\}.\par\medskip
\ifdefined\hypertarget\hypertarget{sourcepointer-4}{}\fi
\noindent\textit{Source reference.} Original lemma 030V: sources/stacks-chow/tags.txt, source lines 3962--3962. Cross-chapter standard prerequisite invoked once at proof lines 457-459 to attach conditions (3)-(6); the cycle (1)-(2)-(7) is retained in the counted span.\par\medskip
\begin{lemma}
\label{lemma-geometrically-reduced}
Let $k$ be a field of characteristic $p > 0$. Let $X$ be a scheme over $k$.
The following are equivalent
\begin{enumerate}
\item $X$ is geometrically reduced,
\item $X_{k'}$ is reduced for every field extension $k'/k$,
\item $X_{k'}$ is reduced for every finite purely inseparable field extension
$k'/k$,
\item $X_{k^{1/p}}$ is reduced,
\item $X_{k^{perf}}$ is reduced,
\item $X_{\bar k}$ is reduced,
\item for every affine open $U \subset X$ the ring $\mathcal{O}_X(U)$
is geometrically reduced (see
Algebra, Definition \ref{algebra-definition-geometrically-reduced}).
\end{enumerate}
\end{lemma}

\begin{proof}
Assume (1). Then for every field extension $k'/k$ and
every point $x' \in X_{k'}$ the local ring of $X_{k'}$ at $x'$
is reduced. In other words $X_{k'}$ is reduced. Hence (2).

\medskip\noindent
Assume (2). Let $U \subset X$ be an affine open. Then for
every field extension $k'/k$ the scheme $X_{k'}$ is reduced, hence
$U_{k'} = \Spec(\mathcal{O}(U)\otimes_k k')$ is reduced,
hence $\mathcal{O}(U)\otimes_k k'$ is reduced (see Properties,
Section \ref{properties-section-integral}). In other words
$\mathcal{O}(U)$ is geometrically reduced, so (7) holds.

\medskip\noindent
Assume (7). For any field extension $k'/k$ the base
change $X_{k'}$ is gotten by gluing the spectra of the
rings $\mathcal{O}_X(U) \otimes_k k'$ where $U$ is affine open
in $X$ (see Schemes, Section \ref{schemes-section-fibre-products}).
Hence $X_{k'}$ is reduced. So (1) holds.

\medskip\noindent
This proves that (1), (2), and (7) are equivalent. These are equivalent
to (3), (4), (5), and (6) because we can apply Algebra, Lemma
\ref{algebra-lemma-geometrically-reduced-finite-purely-inseparable-extension}
to $\mathcal{O}_X(U)$ for $U \subset X$ affine open.
\end{proof}

\end{document}
